\subsection{Coordonnées angulaires, normalisation du tour et unités d'action}
\label{sec:ii-angulos}

La grandeur d’action et la coordonnée angulaire appartiennent à des espaces
distincts. La section \(\hbar_{\rm ret}\) occupe la droite d’action
\(\mathcal L_S\) ; \(\eta_{\rm ret}\) est son coefficient dans la base
déclarée. La représentation angulaire agit sur la paire adimensionnelle
\((\alpha,\eta_{\rm ret})\), après sa construction :
\begin{equation}
 \mathcal P(\alpha,\eta_{\rm ret})
 =\bigl([1000\alpha]_{360},[2(\eta_{\rm ret}+\alpha)]_{360}\bigr).
 \label{eq:ii-par-angular}
\end{equation}
La complétion nonadique détermine le facteur 1000 ; la clôture
bidirectionnelle détermine le facteur deux. La coordinatisation circulaire de 360
unités rend visibles les partitions
\(12\cdot30=4\cdot90=3\cdot120\). Sur la branche positive de
l’article, les représentants sont
\[
 A_{\deg}=1000\alpha,\qquad
 C^*_{\deg}=2(\eta_{\rm ret}+\alpha).
\]

\begin{proposition}[Décalage de la représentation en base mille]
\label{prop:ii-desplazamiento-angular}
Soit \(\alpha=\sum_{n\geq1}b_n1000^{-n}\), avec
\(b_n\in\{0,\ldots,999\}\), le développement compatible déjà engendré.
La coordonnée relevée \(A_{\deg}=1000\alpha\) satisfait
\[
 A_{\deg}=b_1+\sum_{n\geq1}b_{n+1}1000^{-n}.
\]
La représentation conserve la suite des blocs à partir du deuxième ;
le bloc initial passe à la partie entière.
\end{proposition}
\begin{proof}
La série converge absolument. Multiplier chaque terme par 1000
et séparer le terme initial fournit l’identité. L’opération
agit sur le développement construit, avec ses blocs et sa mémoire ;
la représentation circulaire ultérieure prend sa classe modulo 360.
\end{proof}

\begin{definition}[Coordinatisations de la coordonnée angulaire]
La coordinatisation en degrés exprime un tour complet au moyen de 360 unités.
La coordinatisation en radians l’exprime au moyen de \(2\pi\). Le changement de
coordonnées et ses inverses sont
\[
 x=\frac{\pi}{180}A_{\deg},\qquad
 y=\frac{\pi}{180}C^*_{\deg},\qquad
 A_{\deg}=\frac{180}{\pi}x,\qquad
 C^*_{\deg}=\frac{180}{\pi}y.
\]
Sur un parcours dont le nombre de tours est conservé, on utilise les
coordonnées relevées ; le quotient circulaire exprime sa classe modulo
le tour correspondant.
\end{definition}

La définition caractérise les degrés et les radians par leur relation exacte
avec le tour. Le registre retient en outre l’orientation, les tours et la
mémoire de transport. La représentation de sa classe angulaire conserve
la phase, tandis que le relèvement distingue des parcours ayant la même
phase finale. La lecture physique d’action nécessite de préciser laquelle de
ces coordonnées intervient dans l’intégrale.

\begin{proposition}[Covariance du lecteur d'action]
Pour \(h_a=2\pi\hbar_a\), l'action d'une coordonnée relevée
satisfait
\begin{equation}
 \mathcal S_a(\theta)
 =\hbar_a\theta_{\rm rad}
 =h_a\frac{\theta_{\deg}}{360}.
 \label{eq:ii-accion-cartas}
\end{equation}
Un tour positif exprime \(h_a\).
\end{proposition}
\begin{proof}
Substituer \(\theta_{\rm rad}=\pi\theta_{\deg}/180\) donne
\(\hbar_a\pi\theta_{\deg}/180
=(2\pi\hbar_a)\theta_{\deg}/360\). Pour un tour,
\(\theta_{\deg}=360\), d'où \(h_a\).
\end{proof}

La périodicité d'une représentation est déclarée séparément. Si la
représentation transportée satisfait \(U(2\pi)=-I\), sa restitution
opératoire se produit en \(4\pi\). Cette propriété du relèvement
coexiste avec la normalisation par tour de
\eqref{eq:ii-accion-cartas}. Chacune conserve son objet : grandeur
intégrée, phase circulaire et état représenté.

\subsection{Covariance des canaux et de la fonctionnelle d'incidence}

Les canaux s'écrivent de manière équivalente
\begin{equation}
 q_\pm=\exp[-(x\pm y)]
 =\exp\!\left[-\frac{\pi}{180}(A_{\deg}\pm C^*_{\deg})\right].
\label{eq:ii-canales}
\end{equation}
\begin{proposition}[Récupération des coordonnées paire et impaire]
\label{prop:ii-inversion-canales-angulares}
Dans la coordinatisation réelle relevée, les canaux positifs déterminent
de manière univoque les deux coordonnées :
\[
 A_{\deg}=-\frac{90}{\pi}\log(q_+q_-),\qquad
 C^*_{\deg}=\frac{90}{\pi}\log\frac{q_-}{q_+}.
\]
L'échange \((q_+,q_-)\mapsto(q_-,q_+)\) conserve
\(A_{\deg}\) et inverse \(C^*_{\deg}\).
\end{proposition}
\begin{proof}
La somme et la différence des logarithmes de
\eqref{eq:ii-canales} sont respectivement
\(-\pi A_{\deg}/90\) et \(-\pi C^*_{\deg}/90\).
Le logarithme réel est univoque sur les canaux positifs.
Le produit est symétrique et le quotient s'inverse lors de leur échange.
\end{proof}
Dans le registre à douze positions, l'involution antipodale
\(r\mapsto r+6\pmod {12}\) forme six paires. La distribution
uniforme de la coordonnée impaire sur ces paires est
\(C^*_{\deg}/6\) : sommer ses six contributions restitue
\(C^*_{\deg}\). Cette coordonnée sectorielle se distingue de
la constante de Planck \(h\), qui appartient à la droite d'action.

La conversion est appliquée une seule fois à chaque argument. La grandeur
\(\alpha\) demeure adimensionnelle et la grandeur \(\hbar\)
appartient toujours à la droite d'action ; leurs coordonnées antérieures
interviennent dans l'application \(\mathcal P\).

La fonctionnelle de Barbero--Immirzi contient, outre les canaux, une
composante bilinéaire dans les représentants angulaires. Son transport est
\begin{equation}
 \frac{\pi}{180}A_{\deg}C^*_{\deg}
 =\frac{180}{\pi}xy.
 \label{eq:ii-barbero-cartas}
\end{equation}
L'identité s'obtient en remplaçant les deux représentants en degrés
par leurs expressions en radians. La formule complète en
coordonnées analytiques est donc
\[
 \gamma_{\HMT}
 =\frac{180}{\pi}xy
 +12\bigl[S_{90}(e^{-x+y})-S_{90}(e^{-x-y})\bigr]
 -\bigl[S_{120}(e^{-x+y})-S_{120}(e^{-x-y})\bigr].
\]
L’équivalence des deux coordinatisations transporte la composante bilinéaire
et les fonctions exponentielles simultanément. Cette transformation
explicite le rôle précis du degré dans la construction et permet
de reproduire l’évaluation sans ambiguïté d’unités angulaires.
